\newpage
\phantomsection
\label{prf:power-function-continuous}

\begin{remark*}[Return]
\hyperref[lem:power-function-continuous]{Return to Theorem}
\end{remark*}

\begin{theorem*}[Power Function Is Continuous]
Let $f : \mathbb{R} \to \mathbb{R}$ be defined by $f(x) = x^n$ for some $n \in \mathbb{N}$. Then $f$ is a continuous function.
\end{theorem*}

\begin{proof}
\textbf{Professional Standard Proof.}~\\
The proof proceeds by induction on $n$. For the base case $n = 1$, let $x_0 \in \mathbb{R}$ and $\varepsilon > 0$. Setting $\delta := \varepsilon$, one has $|x - x_0| < \delta \Rightarrow |f(x) - f(x_0)| = |x - x_0| < \varepsilon$, establishing continuity of $f(x) = x$. For the inductive step, assume $g(x) = x^n$ is continuous. Then $x^{n+1} = x^n \cdot x = g(x) \cdot \mathrm{id}(x)$ where both $g$ and $\mathrm{id}$ are continuous. By the Product Rule for Continuous Functions, their pointwise product $f(x) = x^{n+1}$ is continuous. By induction, $f(x) = x^n$ is continuous for all $n \in \mathbb{N}$.
\end{proof}

\begin{proof}
\textbf{Detailed Learning Proof.}~\\

\textbf{Step 1.} Establish the base case $n = 1$. Let $x_0 \in \mathbb{R}$ and let $\varepsilon > 0$. Set $\delta := \varepsilon$. For any $x \in \mathbb{R}$ satisfying $|x - x_0| < \delta$, one has $|f(x) - f(x_0)| = |x - x_0| < \delta = \varepsilon$. This verifies the $\varepsilon$-$\delta$ definition of continuity at $x_0$, so $f(x) = x$ is continuous.

\textbf{Step 2.} Formulate the inductive hypothesis. Assume that for some $n \in \mathbb{N}$, the function $g(x) = x^n$ is continuous on $\mathbb{R}$.

\textbf{Step 3.} Prove the inductive step. Consider $f(x) = x^{n+1}$. Rewrite this as $x^{n+1} = x^n \cdot x = g(x) \cdot \mathrm{id}(x)$, where $g(x) = x^n$ is continuous by the inductive hypothesis and $\mathrm{id}(x) = x$ is continuous by Step 1.

\textbf{Step 4.} Apply the Product Rule for Continuous Functions. Since both $g$ and $\mathrm{id}$ are continuous functions, their pointwise product $f(x) = g(x) \cdot \mathrm{id}(x) = x^{n+1}$ is continuous.

\textbf{Step 5.} Conclude by mathematical induction. The base case and inductive step establish that $f(x) = x^n$ is continuous for all $n \in \mathbb{N}$.
\end{proof}

\begin{remark*}[Proof structure]
This is a proof by mathematical induction on the exponent $n$. The base case uses a direct $\varepsilon$-$\delta$ argument with the simple choice $\delta := \varepsilon$. The inductive step avoids complicated algebraic manipulations by recognizing $x^{n+1}$ as the product of two functions already known to be continuous, then invoking the product rule rather than returning to the definition of continuity.
\end{remark*}

\begin{remark*}[Dependencies]~\\
\begin{itemize}
  \item \hyperref[def:continuous-at-point]{Definition of continuity at a point}
  \item \hyperref[thm:product-rule-continuous-functions]{Product Rule for Continuous Functions}
\end{itemize}
\end{remark*}
